%% chapters/00-abstract-en.tex
%% Abstract in English (~500–800 words, mirrors Indonesian abstract structure)
%% Pedoman ITB Subbab~II.2: 4 paragraphs, no references, keywords ≤ 7 phrases
%% Phase H — rewritten to reflect two-track narrative per outline outline Abstract

\begin{abstracten}

Rolling element bearings are the most failure-prone components in industrial
rotating machinery, accounting for approximately 40--50\% of total machine
failures.
The transformation of Indonesian manufacturing toward Industry~4.0 drives
adoption of condition-based predictive maintenance in place of conventional
corrective and preventive strategies, owing to its ability to detect
incipient faults and forecast remaining useful life (RUL) accurately.
Two research gaps motivate this dissertation: first, explainable AI (XAI)
methods for diagnostic deep learning still operate at the aggregated feature
level, without reaching raw vibration signals at 2,048 points; second,
deep learning-based RUL estimation models have not been proven to
internalize bearing characteristic frequencies (BPFO, BPFI, BSF, FTF)
in their latent representations in a statistically rigorous manner.

The first research track benchmarks three algorithm groups for bearing
fault classification on the CWRU dataset (10 fault classes, 2,048-point
signals, 54/13/33\% temporal split): kernel-based models (SVM with RBF
kernel and one-vs-rest Logistic Regression), tree-based algorithms (DT,
RF, and XGBoost), and the Wide and Deep Convolutional Neural Network
(WDCNN) architecture.
WDCNN achieves 99.87\% accuracy (749 of 750 test samples, macro F1 = 0.997),
surpassing SVM-RBF (96.4\%) and XGBoost (97.1\%).
SHAP DeepExplainer attributions over raw signals are aggregated into
per-class Fault Signature Maps (FSM), yielding a split-half stability
coefficient of 0.940 and a discriminability index of 0.216.
An ablation study across three WDCNN variants produces a finding unreported
in the bearing classification literature: removing BatchNorm increases FSM
discriminability by 133\% (from 0.216 to 0.735), at the cost of a 3.60
percentage-point accuracy drop.

The second track develops and evaluates three RUL estimation backbone
architectures: Mamba-xLSTM-Net (898K parameters, combining a selective
state-space model with matrix-memory recurrent LSTM), N-BEATS-xLSTM-RUL
(459K parameters, exploiting structural priors on degradation curve shape),
and SparseGate-TCN-RUL (249K parameters, a lightweight model suited for
edge deployment).
All three are trained under a shared protocol (75 epochs, bf16 precision,
batch size 512, window 32, three random seeds) on PHM2012, XJTU-SY, and
IMS.
The winning model differs by dataset: SparseGate-TCN achieves
RMSE $0.226 \pm 0.030$ on PHM2012, \textcolor{blue}{Mamba-xLSTM-Net achieves
RMSE $0.213 \pm 0.056$ on XJTU-SY}, and Mamba-xLSTM achieves
$0.407 \pm 0.040$ on IMS.
After training, a Top-$k$ Sparse Autoencoder (SAE, $d_{\rm lat} = 1{,}024$,
$k = 51$) is trained post-hoc on the frozen hidden states of each backbone.
Mapping to characteristic frequencies via Pearson correlation and Hilbert
envelope spectra shows that 2.3\% of SAE features exhibit moderate
correlation ($|r| \geq 0.30$) with BPFI on PHM2012 ($p < 0.001$,
$r_{\rm max} = 0.507$), \textcolor{blue}{while on XJTU-SY the hit-rate spreads across BPFI, BSF, and BPFO
(1.5\%, 1.6\%, 0.7\%; $p < 0.001$, $r_{\rm max} = 0.468$).}
Two negative controls, namely an untrained Xavier-initialized model and
Gaussian-noise hidden states, confirm that this correspondence is an
emergent property of trained representations, not an artifact of the
architecture or the SAE procedure.

This dissertation contributes five novelties: N1 (FSM, the first
signal-level XAI for WDCNN in bearing diagnostics), N2 (the BatchNorm
discriminability trade-off, a finding unreported in the bearing
classification literature), N3 (SAEBearing, the first adaptation of
mechanistic interpretability to bearing prognostics), N4 (cross-architecture
universality of BPFx correspondence across three architectures and four
datasets, verified with bootstrap CI and permutation tests), and N5 (a
conceptual three-tier PdM framework: Edge IoT via kernel-based models,
Edge Server via WDCNN+FSM, and Cloud/GPU via RUL backbone and SAE-BPFx).
Taken together, these contributions form the first statistically verified
bridge between data-driven deep learning and classical bearing vibration
theory: the diagnostic track connects raw signals to FSM, while the
prognostic track connects latent representations to BPFx characteristic
frequencies.

\bigskip\noindent
\textbf{Keywords:} predictive maintenance, bearing diagnostics, Fault
Signature Maps, remaining useful life, Sparse Autoencoder, mechanistic
interpretability, bearing characteristic frequencies

\end{abstracten}
